\section{Deskripsi Penugasan}\label{sec:Pendahuluan}

Ice Sliding Puzzle adalah permainan pencarian lintasan pada papan dua
dimensi. Pemain berada pada suatu petak awal dan harus mencapai petak
tujuan. Setiap kali pemain memilih arah gerak, pemain akan terus
meluncur pada arah tersebut sampai bertemu penghalang. Selama
meluncur, pemain dapat melewati beberapa jenis petak, seperti petak
biasa, checkpoint, lava, dan dinding. Jika pemain keluar papan atau
melewati lava, langkah tersebut dianggap tidak valid.

Tugas kecil ini meminta dibuatnya program yang dapat menyelesaikan
Ice Sliding Puzzle dengan pendekatan pencarian graf. Program menerima
konfigurasi papan dari berkas teks, membangun ruang status permainan
secara implisit, lalu mencari urutan gerakan yang membawa pemain dari
posisi awal menuju tujuan. Jika terdapat checkpoint, checkpoint harus
dilewati sesuai urutan sebelum permainan dianggap selesai.

\section{Spesifikasi Tugas}\label{sec:Spesifikasi Tugas}

\subsection{Spesifikasi Wajib}\label{sub:Spesifikasi Wajib}

Program wajib memenuhi kebutuhan berikut.
\begin{itemize}[label=\textbullet]
    \item Membaca masukan papan dari berkas \texttt{.txt}.
    \item Memvalidasi format papan, posisi awal, posisi tujuan,
        checkpoint, dan matriks biaya.
    \item Menyelesaikan persoalan menggunakan algoritma wajib UCS,
        GBFS, dan A*.
    \item Menampilkan solusi berupa urutan gerakan, total biaya,
        jumlah iterasi, waktu eksekusi, dan visualisasi langkah solusi.
    \item Menyediakan penyimpanan hasil solusi ke berkas
        \texttt{.txt}.
    \item Menyediakan antarmuka CLI untuk menjalankan solver dari terminal.
\end{itemize}

Format umum berkas masukan adalah sebagai berikut.
\begin{tcolorbox}[colback=white, colframe=black, sharp corners, boxrule=0.5pt]
\begin{verbatim}
N M
<N baris grid>
<N x M nilai cost>
\end{verbatim}
\end{tcolorbox}

Karakter pada grid memiliki arti sebagai berikut.
\begin{table}[H]
    \centering
    \begin{tabular}{|c|L{9cm}|}
        \hline
        \textbf{Karakter} & \textbf{Keterangan} \\
        \hline
        \texttt{X} & Dinding. Pemain berhenti sebelum memasuki petak ini. \\
        \hline
        \texttt{*} & Petak biasa atau ice tile. \\
        \hline
        \texttt{L} & Lava. Langkah yang melewati petak ini tidak valid. \\
        \hline
        \texttt{Z} & Posisi awal pemain. \\
        \hline
        \texttt{O} & Posisi tujuan. \\
        \hline
        \texttt{0}--\texttt{9} & Checkpoint yang harus dilewati berurutan. \\
        \hline
    \end{tabular}
    \caption{Daftar karakter pada papan}
\end{table}

\subsection{Spesifikasi Bonus}\label{sub:Spesifikasi Bonus}

Pada implementasi ini juga dibuat beberapa fitur bonus. Pertama, GUI
web dibuat agar pengguna dapat mengunggah testcase, memilih algoritma,
memilih heuristic untuk GBFS dan A*, melihat animasi solusi per
langkah, serta melihat trace graf pencarian yang terbentuk selama
solver berjalan. GUI juga menyediakan tombol untuk menyimpan hasil
solusi ke berkas \texttt{.txt}. Kedua, program menambahkan algoritma pathfinding
alternatif BFS dan DFS. Ketiga, selain heuristic utama, program
menyediakan beberapa alternatif heuristic untuk GBFS dan A*.

\section{Dasar Teori}\label{sec:Dasar Teori}

\subsection{Representasi Graf}\label{sub:Representasi Graf}

Persoalan Ice Sliding Puzzle dapat dipandang sebagai persoalan
pencarian lintasan pada graf berbobot. Setiap simpul menyatakan suatu
status permainan, yaitu posisi pemain dan checkpoint berikutnya yang
harus dicapai. Setiap sisi menyatakan satu aksi slide ke salah satu
arah. Bobot sisi adalah total biaya petak yang dilewati selama slide
tersebut.

Graf tidak perlu dibangun seluruhnya sejak awal. Program cukup
membangkitkan tetangga dari simpul yang sedang diekspansi. Pendekatan
ini disebut pembangkitan graf implisit dan cocok untuk ruang status
yang dapat membesar secara cepat.

\subsection{Uniform Cost Search}\label{sub:UCS}

Uniform Cost Search (UCS) adalah algoritma pencarian yang selalu
mengekspansi simpul dengan biaya lintasan terkecil dari simpul awal.
Jika semua biaya sisi tidak negatif, UCS menjamin solusi optimal
berdasarkan total biaya. Prioritas simpul pada UCS didefinisikan
sebagai berikut.
\[
    f(n) = g(n)
\]
dengan $g(n)$ adalah biaya lintasan dari simpul awal menuju simpul
$n$.

\subsection{Greedy Best First Search}\label{sub:GBFS}

Greedy Best First Search (GBFS) memilih simpul berdasarkan nilai
heuristic yang memperkirakan jarak menuju target. Algoritma ini
cenderung cepat menuju daerah yang dianggap menjanjikan, tetapi tidak
menjamin solusi optimal karena biaya lintasan yang sudah ditempuh
tidak menjadi bagian dari prioritas utama. Prioritas GBFS adalah:
\[
    f(n) = h(n)
\]

\subsection{A*}\label{sub:AStar}

A* menggabungkan biaya aktual yang sudah ditempuh dengan estimasi
biaya menuju target. Prioritas simpul pada A* adalah:
\[
    f(n) = g(n) + h(n)
\]
Jika heuristic yang digunakan admissible dan konsisten, A* dapat
menemukan solusi optimal dengan jumlah ekspansi yang umumnya lebih
sedikit dibanding UCS.

\subsection{Breadth First Search}\label{sub:BFS}

Breadth First Search (BFS) adalah algoritma pencarian tidak berbobot
yang mengekspansi simpul berdasarkan kedalaman. Pada implementasi ini,
setiap aksi slide dianggap sebagai satu langkah, sehingga BFS optimal
terhadap jumlah slide, bukan terhadap total cost traversal.
\[
    f(n) = depth(n)
\]

\subsection{Depth First Search}\label{sub:DFS}

Depth First Search (DFS) menggunakan stack untuk menelusuri cabang
sedalam mungkin sebelum mencoba cabang lain. DFS dapat menemukan
solusi pada ruang status finite, tetapi tidak menjamin solusi dengan
jumlah slide paling sedikit maupun cost paling kecil.

\subsection{Heuristic}\label{sub:Heuristic}

Program menyediakan tiga mode heuristic.
\begin{enumerate}
    \item \texttt{H1}: heuristic utama berupa jarak Manhattan dari
        posisi sekarang ke target berikutnya.
    \item \texttt{H2}: alternatif heuristic berupa jarak Manhattan ke
        target berikutnya dikalikan biaya petak minimum yang dapat dilewati.
    \item \texttt{H3}: alternatif heuristic berupa total jarak
        Manhattan menuju checkpoint yang
        tersisa secara berurutan, lalu menuju goal.
\end{enumerate}

Pada permainan slide, jarak Manhattan bukan jarak langkah sebenarnya
karena satu gerakan dapat melewati beberapa petak. Namun, heuristic
ini tetap berguna sebagai estimasi arah pencarian yang sederhana dan
murah dihitung.

Untuk A*, admissibility heuristic pada permainan ini tidak selalu
dapat dijamin secara umum. Satu aksi slide dapat melewati beberapa
petak sekaligus, sehingga jarak Manhattan dalam satuan petak dapat
lebih besar daripada biaya aksi minimum jika terdapat tile berbiaya
rendah atau jika satu slide mendekatkan posisi melewati banyak petak.
\texttt{H1} dan \texttt{H2} tetap dipakai karena murah dihitung dan
memberi arah pencarian yang baik, tetapi optimalitas A* secara teoritis
hanya terjamin pada testcase ketika nilai heuristic tidak melebihi
biaya sisa sebenarnya. \texttt{H3} lebih informatif untuk checkpoint
berurutan karena menjumlahkan estimasi ke target-target yang tersisa,
namun juga tidak dijamin admissible pada seluruh konfigurasi papan
karena tidak memodelkan titik berhenti akibat dinding dan efek slide.
